\documentclass[12pt, a4paper]{article}

\usepackage[french]{babel}
\usepackage[utf8]{inputenc}
\usepackage[T1]{fontenc}
\usepackage{tgadventor}
\renewcommand*\familydefault{\sfdefault}
\usepackage{amsmath, amssymb}
\usepackage[top=2cm, bottom=2cm, left=2cm, right=2cm]{geometry}
\usepackage{enumitem}
\usepackage{multicol}
\usepackage[table]{xcolor}
\usepackage{mdframed}
\usepackage{verbatim}
\usepackage{tabularx}
\usepackage{tkz-euclide}

\definecolor{vlmblue}{RGB}{25, 90, 160}
\definecolor{vlmgray}{RGB}{245, 245, 245}
\definecolor{vlmgreen}{RGB}{0, 130, 80}

% ── Titre de l'exercice dans un encadré ──
\newcommand{\titrexercice}[2]{%
  \begin{mdframed}[
    backgroundcolor=vlmgray,
    linecolor=black,
    linewidth=1.5pt,
    innerleftmargin=10pt,
    innerrightmargin=10pt,
    innertopmargin=6pt,
    innerbottommargin=6pt,
    skipabove=1em,
    skipbelow=0.5em
  ]
  \textbf{Exercice #1}%
  \ifx&#2&\else\hfill\textbf{#2}\fi
  \end{mdframed}
}

% ── Corps de l'exercice sans cadre ──
\newenvironment{exercice}[2]{%
  \titrexercice{#1}{#2}
  \begin{list}{}{%
    \setlength{\leftmargin}{0pt}
    \setlength{\rightmargin}{0pt}
    \setlength{\listparindent}{0pt}
    \setlength{\itemindent}{0pt}
    \setlength{\parsep}{0pt}
  }\item[]\noindent\ignorespaces
}{%
  \end{list}
  \vspace{1em}
}

% ── Solution ──
\ifdefined\AVECSOLUTIONS
  \newenvironment{solution}[1]{%
    \vspace{0.3em}
    \textbf{\textcolor{vlmgreen}{Solution #1}}
    \vspace{0.3em}\par
  }{%
    \vspace{1em}
  }
\else
  \newenvironment{solution}[1]{\comment}{\endcomment}
\fi

\pagestyle{empty}

\begin{document}
\begin{exercice}{EX-CH02-011}{Additions et soustractions de nombres relatifs}
{\small
\bigskip
\def\arraystretch{1.7}
\begin{tabularx}{\linewidth}{*{4}{X}}
\textbf{A.} & \textbf{B.} & \textbf{C.} & \textbf{D.} \\
$(+8)+(-3)= \dotfill$   & $(-5)+(+9)= \dotfill$   & $(+10)+(-14)= \dotfill$ & $(-3)-(+9)= \dotfill$   \\
$(+6)-(+12)= \dotfill$  & $(+14)+(-19)= \dotfill$ & $(-10)+(-12)= \dotfill$ & $(+8)-(-8)= \dotfill$   \\
$(+4)+(+11)= \dotfill$  & $(-17)-(-15)= \dotfill$ & $(+9)+(-9)= \dotfill$   & $(+24)-(+14)= \dotfill$ \\
$(-24)-(+12)= \dotfill$ & $(+9)+(+13)= \dotfill$  & $(-12)+(-8)= \dotfill$  & $(-14)+(-7)= \dotfill$  \\
$(+15)-(-7)= \dotfill$  & $(+10)-(-9)= \dotfill$  & $(+14)+(-5)= \dotfill$  & $(-15)-(-17)= \dotfill$ \\
$(+14)+(-8)= \dotfill$  & $(+12)-(-1)= \dotfill$  & $(-5)-(-14)= \dotfill$  & $(+20)-(+18)= \dotfill$ \\
\end{tabularx}
}
\end{exercice}
\begin{solution}{EX-CH02-011}

\def\arraystretch{1.7}
\begin{tabularx}{\linewidth}{*{4}{X}}
\textbf{A.} & \textbf{B.} & \textbf{C.} & \textbf{D.} \\
$(+8)+(-3)=+5$   & $(-5)+(+9)=+4$   & $(+10)+(-14)=-4$ & $(-3)-(+9)=-12$  \\
$(+6)-(+12)=-6$  & $(+14)+(-19)=-5$ & $(-10)+(-12)=-22$& $(+8)-(-8)=+16$  \\
$(+4)+(+11)=+15$ & $(-17)-(-15)=-2$ & $(+9)+(-9)=0$    & $(+24)-(+14)=+10$\\
$(-24)-(+12)=-36$& $(+9)+(+13)=+22$ & $(-12)+(-8)=-20$ & $(-14)+(-7)=-21$ \\
$(+15)-(-7)=+22$ & $(+10)-(-9)=+19$ & $(+14)+(-5)=+9$  & $(-15)-(-17)=+2$ \\
$(+14)+(-8)=+6$  & $(+12)-(-1)=+13$ & $(-5)-(-14)=+9$  & $(+20)-(+18)=+2$ \\
\end{tabularx}

\end{solution}
\end{document}